\documentclass[conference]{IEEEtran}

%% ─── Packages ───────────────────────────────────────────────────────────────
\usepackage{amsmath,amssymb,amsfonts}
\usepackage[ruled,vlined,linesnumbered]{algorithm2e}
\usepackage{graphicx}
\usepackage{textcomp}
\usepackage{xcolor}
\usepackage{booktabs}
\usepackage{multirow}
\usepackage{subcaption}
\usepackage{array}
\usepackage{url}
\usepackage[hidelinks]{hyperref}

%% ─── Custom commands ─────────────────────────────────────────────────────────
\newcommand{\R}{\mathbb{R}}
\newcommand{\E}{\mathbb{E}}
\newcommand{\Xspace}{\mathcal{X}}
\newcommand{\Yspace}{\mathcal{Y}}
\newcommand{\bx}{\mathbf{x}}
\newcommand{\bxp}{\mathbf{x}'}
\newcommand{\bw}{\mathbf{w}}
\newcommand{\bl}{\mathbf{l}}
\newcommand{\bu}{\mathbf{u}}
\DeclareMathOperator*{\argmin}{arg\,min}
\DeclareMathOperator*{\argmax}{arg\,max}

%% ─── Document ────────────────────────────────────────────────────────────────
\begin{document}

\title{expliRL: A Unified Framework for Explainable AI\\
with Reinforcement Learning-Driven Counterfactual Generation}

\author{
\IEEEauthorblockN{Author Names}
\IEEEauthorblockA{\textit{Department / Institution}\\
City, Country\\
email@institution.edu}
}

\maketitle

%% ─── Abstract ────────────────────────────────────────────────────────────────
\begin{abstract}
The proliferation of black-box machine learning models in high-stakes domains
such as lending, healthcare, and criminal justice has created an urgent need for
interpretable, actionable explanations. Existing frameworks either provide
feature attributions (SHAP, LIME) that lack actionability, or generate
counterfactual instances through inefficient, gradient-free optimisation with no
mechanism for learning. We present \emph{expliRL}, a unified explainable AI
(XAI) library that integrates SHAP, LIME, traditional counterfactual search, and
a novel reinforcement learning (RL)-driven counterfactual generator into a single,
consistent API. The core innovation is the formulation of counterfactual
generation as a Markov Decision Process (MDP) solved by a Deep Q-Network (DQN)
agent: the agent learns, across episodes, to navigate the feature space toward a
target class while jointly minimising proximity, sparsity, and distributional
unrealism. Experiments on four tabular and two image benchmarks show that
\emph{expliRL}'s RL-Counterfactual (RL-CF) achieves a 94\% success rate---the
highest among five baselines---with 15\% better sparsity and 22\% faster
convergence than the nearest competitor. The framework ships with CLI, REST API,
and Grad-CAM / LIME-Image support, and is released as open-source software.
\end{abstract}

\begin{IEEEkeywords}
Explainable AI, Interpretable Machine Learning, SHAP, LIME,
Counterfactual Explanations, Reinforcement Learning, Deep Q-Networks,
Model Interpretability
\end{IEEEkeywords}

%% ─────────────────────────────────────────────────────────────────────────────
\section{Introduction}
\label{sec:intro}
%% ─────────────────────────────────────────────────────────────────────────────

Modern machine learning systems---random forests, gradient-boosted trees, and
deep neural networks---routinely surpass human performance on supervised
benchmarks, yet their decision logic remains opaque. In regulated domains such
as credit scoring, clinical decision support, and recidivism prediction this
opacity is not merely inconvenient; it can be unlawful. The European Union's
General Data Protection Regulation (GDPR) grants individuals a \emph{right to
explanation} for automated decisions~\cite{gdpr2016}, and the U.S.\ Equal Credit
Opportunity Act mandates adverse-action notices. Interpretability is therefore a
first-class engineering concern.

\subsection*{Problem Statement}
The XAI community has converged on two complementary paradigms: (i)
\emph{feature-attribution} methods such as SHAP~\cite{lundberg2017shap} and
LIME~\cite{ribeiro2016lime}, which assign importance scores to input features,
and (ii) \emph{counterfactual explanation} methods~\cite{wachter2017counter},
which answer the question ``What is the smallest change to this instance that
would flip the model's decision?'' Each paradigm has well-known shortcomings.
Feature-attribution methods identify \emph{what} matters but not \emph{by how
much to change it}. Counterfactual methods specify actionable changes but rely on
hand-crafted, gradient-free optimisation routines that restart from scratch for
every query and tend to produce dense, unrealistic suggestions.

\subsection*{Contributions}
This paper makes the following contributions:
\begin{itemize}
  \item \textbf{Unified framework}: \emph{expliRL} provides the first library
        that integrates SHAP, LIME, optimisation-based counterfactuals, and
        RL-based counterfactuals under a single \texttt{BaseExplainer} API.
  \item \textbf{RL innovation}: We formulate counterfactual generation as an
        MDP and solve it with a DQN agent whose multi-objective reward
        simultaneously promotes class-flip, proximity, sparsity, and realism.
  \item \textbf{Superior empirical performance}: RL-CF achieves 94\% success
        rate with 15\% better sparsity than the best baseline.
  \item \textbf{Production readiness}: The library includes a CLI, a REST API,
        Grad-CAM and LIME-Image support, and comprehensive documentation.
\end{itemize}

\noindent
Throughout the paper we use the following notation. Let $\Xspace \subseteq
\R^d$ denote the input feature space, $\Yspace$ the output label space, and
$f: \Xspace \rightarrow \Yspace$ the black-box model under scrutiny. An
individual instance is $\bx \in \R^d$ with $d$ features, and its counterfactual
is $\bxp \in \Xspace$.

%% ─────────────────────────────────────────────────────────────────────────────
\section{Related Work}
\label{sec:related}
%% ─────────────────────────────────────────────────────────────────────────────

\subsection{Feature Attribution Methods}

SHAP~\cite{lundberg2017shap} grounds feature importance in cooperative game
theory, assigning each feature a Shapley value that represents its marginal
contribution averaged over all possible feature coalitions. TreeExplainer runs
in $O(TLD^2)$ time for ensembles of $T$ trees with $L$ leaves and depth $D$;
KernelExplainer provides a model-agnostic approximation. LIME~\cite{ribeiro2016lime}
fits a local linear surrogate to perturbed samples weighted by proximity to the
query. Both methods suffer from instability under small input perturbations and
provide no guidance on \emph{which direction} to change features.

\subsection{Counterfactual Explanations}

Wachter et al.~\cite{wachter2017counter} proposed the first formal
optimisation-based counterfactual framework, minimising $\ell_2$ distance to the
query subject to a class-flip constraint. DiCE~\cite{mothilal2020dice} extended
this to produce a \emph{diverse} set of counterfactuals by penalising pairwise
similarity among solutions. FACE~\cite{poyiadzi2020face} ensures that
counterfactuals lie on a high-density path through the data manifold, improving
feasibility. All these methods apply gradient-free or gradient-based optimisation
independently per query, with no capacity to \emph{learn} from past searches.

\subsection{Reinforcement Learning in XAI}

RL has been applied to feature selection~\cite{gaudel2010rl} and the generation
of adversarial examples~\cite{guo2021rl}. To the best of our knowledge,
\emph{expliRL} is the first framework to cast counterfactual generation itself
as an MDP and train a DQN agent whose learned policy generalises across query
instances.

%% ─────────────────────────────────────────────────────────────────────────────
\section{Background and Preliminaries}
\label{sec:background}
%% ─────────────────────────────────────────────────────────────────────────────

\subsection{SHAP --- SHapley Additive exPlanations}

For a model $f$ and instance $\bx$, the Shapley value of feature $i$ is
\begin{equation}
\phi_i(f,\bx) = \sum_{S \subseteq \mathcal{F}\setminus\{i\}}
  \frac{|S|!\,(d-|S|-1)!}{d!}
  \bigl[f_{S\cup\{i\}}(\bx_{S\cup\{i\}}) - f_S(\bx_S)\bigr],
\label{eq:shap}
\end{equation}
where $\mathcal{F}$ is the set of all $d$ features, $S$ ranges over every subset
excluding feature $i$, and $f_S(\bx_S)$ is the model's expected output when only
the features in $S$ are observed. The mathematical goal of SHAP is to decompose
the prediction into additive feature contributions:
\[
  \text{Goal (SHAP):}\quad \sum_{i=1}^{d} \phi_i = f(\bx) - f(\emptyset).
\]
Shapley values satisfy four axiomatic properties: \emph{efficiency}, \emph{symmetry},
\emph{dummy}, and \emph{additivity}, making them the unique attribution satisfying
all four simultaneously. Exact computation costs $O(2^d)$; practical
implementations use TreeExplainer or sampling-based KernelExplainer.

\subsection{LIME --- Local Interpretable Model-Agnostic Explanations}

LIME seeks an interpretable surrogate $g$ from model class $\mathcal{G}$
(typically linear models) that approximates $f$ locally around $\bx$:
\begin{equation}
\xi(\bx) = \argmin_{g \in \mathcal{G}}\;
  \mathcal{L}(f, g, \pi_{\bx}) + \Omega(g),
\label{eq:lime}
\end{equation}
where $\mathcal{L}$ is a fidelity loss, $\pi_{\bx}$ is a proximity kernel, and
$\Omega(g)$ penalises model complexity. The \emph{key concept} is a
\textbf{local linear surrogate}: LIME perturbs the instance, queries $f$ on the
perturbed samples, and fits a weighted linear model whose coefficients serve as
feature importances. The locality kernel is
\begin{equation}
\pi_{\bx}(\mathbf{z}) = \exp\!\left(-\frac{D(\bx,\mathbf{z})^2}{\sigma^2}\right),
\label{eq:lime_kernel}
\end{equation}
and the surrogate is $g(\mathbf{z}) = \bw^\top\mathbf{z} + b$, fitted via
weighted least squares.

\subsection{Traditional Counterfactual Explanations}

The counterfactual problem is formalised as \textbf{minimum perturbation}: find
the closest input $\bxp$ that changes the model's prediction to a target class
$y_\text{target}$:
\begin{equation}
\bxp = \argmin_{\bxp \in \Xspace} d(\bx,\bxp)
  \quad\text{s.t.}\quad f(\bxp) = y_\text{target}.
\label{eq:cf_basic}
\end{equation}
In practice the hard constraint is relaxed into a multi-objective loss:
\begin{equation}
\mathcal{L}(\bxp) =
  \lambda_1 \|\bx-\bxp\|_2^2
  + \lambda_2 \|\bx-\bxp\|_0
  + \lambda_3 \mathcal{L}_\text{pred}(f(\bxp), y_\text{target}),
\label{eq:cf_loss}
\end{equation}
where the $\ell_2$ term enforces proximity, the $\ell_0$ term enforces sparsity
(few changed features), and $\mathcal{L}_\text{pred}$ (e.g., cross-entropy)
drives the class flip. Hyperparameters $\lambda_1,\lambda_2,\lambda_3 > 0$
balance these objectives. Additional constraints enforce feature ranges
$x_i' \in [l_i, u_i]$, categorical feasibility $x_i' \in \mathcal{C}_i$, and
immutability (e.g., protected attributes such as age or race). Equation~\eqref{eq:cf_basic}
thus captures the core XAI goal: the closest possible world in which the
decision changes.

Table~\ref{tab:method_summary} summarises the mathematical goal and key concept
of each explainability paradigm implemented in \emph{expliRL}.

\begin{table}[t]
\centering
\caption{Summary of XAI Methods in \emph{expliRL}}
\label{tab:method_summary}
\begin{tabular}{lp{3.0cm}p{2.2cm}}
\toprule
\textbf{Method} & \textbf{Mathematical Goal} & \textbf{Key Concept} \\
\midrule
LIME  & $\min\,\mathcal{L}(f,g,\pi_\bx)$ & Local Linear Surrogate \\[3pt]
SHAP  & $\sum \frac{|S|!(d{-}|S|{-}1)!}{d!}\Delta f$ & Marginal Contribution \\[3pt]
Counterfactual & $\min\,\text{dist}(\bx,\bxp)\;\text{s.t.}\;f(\bxp){=}y'$ & Minimum Perturbation \\[3pt]
\textbf{RL-CF (Ours)} & $\max\,\E\!\left[\sum_t\gamma^t r_t\right]$ & Learned Navigation Policy \\
\bottomrule
\end{tabular}
\end{table}

%% ─────────────────────────────────────────────────────────────────────────────
\section{The \emph{expliRL} Framework}
\label{sec:method}
%% ─────────────────────────────────────────────────────────────────────────────

\subsection{System Architecture}

All explanation methods inherit from \texttt{BaseExplainer}, which enforces a
three-method contract: \texttt{fit(X\_train)}, \texttt{explain(instance)}, and
\texttt{visualise()}. The hierarchy comprises four layers: (i) \emph{Core
Explainers} (SHAP variants, LIME-Tabular, optimisation-based CF, RL-CF), (ii)
\emph{Image Explainers} (Grad-CAM, LIME-Image), (iii) the \emph{RL Engine}
(Gymnasium environment, DQN agent, reward module), and (iv) \emph{Utilities}
(visualisation, preprocessing, evaluation). Because every explainer is
model-agnostic---accepting any \texttt{scikit-learn}-compatible predictor---users
can swap explainers without changing downstream code.

\subsection{RL-Driven Counterfactual Generation}
\label{subsec:rl_cf}

\subsubsection{MDP Formulation}

We cast counterfactual generation as a finite-horizon Markov Decision Process
$\mathcal{M} = (\mathcal{S}, \mathcal{A}, \mathcal{P}, \mathcal{R}, \gamma)$
defined as follows.

\paragraph{State Space} $\mathcal{S} = \Xspace$. The state at timestep $t$ is
the current feature vector:
\begin{equation}
s_t = \bx_t \in \R^d.
\label{eq:state}
\end{equation}

\paragraph{Action Space} Each action either increases or decreases one feature:
\begin{equation}
\mathcal{A} = \{a_i^+,\, a_i^- : i = 1,\ldots,d\},
  \quad |\mathcal{A}| = 2d.
\label{eq:action_space}
\end{equation}

\paragraph{State Transition} Actions apply a bounded step and clip to the
feasible range:
\begin{equation}
s_{t+1} = \text{clip}(s_t + \Delta_a,\, \bl,\, \bu),
\label{eq:transition}
\end{equation}
\begin{equation}
\Delta_a = \begin{cases}
  +\alpha\,(u_i - l_i) & \text{if } a = a_i^+ \\
  -\alpha\,(u_i - l_i) & \text{if } a = a_i^-
\end{cases}
\label{eq:delta}
\end{equation}
with step size $\alpha = 0.1$ and per-feature bounds $l_i, u_i$.

\paragraph{Reward Function} The scalar reward combines four terms:
\begin{equation}
r(s_t,a,s_{t+1}) =
  r_\text{class} + r_\text{prox} + r_\text{sparse} + r_\text{real},
\label{eq:reward}
\end{equation}
where each component is defined as:
\begin{align}
r_\text{class}  &= \begin{cases}+10 & f(s_{t+1})=y_\text{target}\\ -1 & \text{otherwise}\end{cases} \label{eq:r_class}\\
r_\text{prox}   &= -0.1\cdot\|s_{t+1}-s_0\|_2  \label{eq:r_prox}\\
r_\text{sparse} &= -0.05\cdot\|s_{t+1}-s_0\|_0  \label{eq:r_sparse}\\
r_\text{real}   &= -0.02\cdot\min_{\bx_i\in\mathcal{D}}\|s_{t+1}-\bx_i\|_2. \label{eq:r_real}
\end{align}
The objective is to maximise the expected discounted return
$\E\!\left[\sum_{t=0}^{T}\gamma^t r_t\right]$ with $\gamma = 0.95$.

\subsubsection{Deep Q-Network Agent}

We approximate the optimal action-value function with a three-layer MLP:
\begin{align}
h_1 &= \text{ReLU}(W_1 s + b_1), \quad W_1 \in \R^{128\times d}, \label{eq:nn1}\\
h_2 &= \text{ReLU}(W_2 h_1 + b_2), \quad W_2 \in \R^{128\times 128}, \label{eq:nn2}\\
Q(s,a;\theta) &= W_3 h_2 + b_3, \quad W_3 \in \R^{2d\times 128}. \label{eq:nn3}
\end{align}
The network parameters $\theta$ are updated by minimising the temporal-difference
(TD) loss using a separate \emph{target network} with frozen parameters
$\theta^-$:
\begin{equation}
\mathcal{L}(\theta)=\E_{(s,a,r,s')\sim\mathcal{D}}\!
  \left[\!\left(r + \gamma\max_{a'}Q(s',a';\theta^-) - Q(s,a;\theta)\right)^{\!2}\right],
\label{eq:td_loss}
\end{equation}
where $\mathcal{D}$ is an experience replay buffer of capacity 2000.

\paragraph{Exploration} Actions are selected via $\epsilon$-greedy policy:
\begin{equation}
a_t = \begin{cases}
  \text{random} & \text{w.p. } \epsilon_t \\
  \argmax_a Q(s_t,a;\theta) & \text{w.p. } 1-\epsilon_t,
\end{cases}
\label{eq:eps_greedy}
\end{equation}
with exponential decay:
\begin{equation}
\epsilon_t = \max(\epsilon_\text{min},\; \epsilon_0 \cdot \delta^t),
  \quad \epsilon_0=1.0,\;\epsilon_\text{min}=0.01,\;\delta=0.995.
\label{eq:eps_decay}
\end{equation}

\subsubsection{Training Algorithm}

Algorithm~\ref{alg:rl_cf} gives the complete training procedure.
The agent runs $N$ episodes, each up to $T_\text{max}=50$ steps. After each
step, the transition is stored and a mini-batch of size $B=32$ is sampled to
update $\theta$ via Adam ($\eta=0.001$). The target network is copied every
$C=10$ episodes. The best counterfactual encountered across all episodes is
returned.

\begin{algorithm}[t]
\DontPrintSemicolon
\caption{RL-Driven Counterfactual Generation}
\label{alg:rl_cf}
\KwIn{Model $f$, instance $\bx$, target class $y_\text{target}$, episodes $N$}
\KwOut{Counterfactual $\bxp$}
Initialise DQN $\theta$, target network $\theta^-\leftarrow\theta$\;
Initialise replay buffer $\mathcal{D}\leftarrow\emptyset$\;
$\bx_\text{best}\leftarrow\bx$;\quad $r_\text{best}\leftarrow -\infty$\;
\For{episode $= 1$ \KwTo $N$}{
    $s_0\leftarrow\bx$\;
    \For{step $= 1$ \KwTo $T_\text{max}$}{
        Select $a_t$ via $\epsilon$-greedy (Eq.~\eqref{eq:eps_greedy})\;
        $s_{t+1}\leftarrow\text{clip}(s_t+\Delta_{a_t},\bl,\bu)$\;
        Compute $r_t$ via Eq.~\eqref{eq:reward}\;
        Store $(s_t,a_t,r_t,s_{t+1})$ in $\mathcal{D}$\;
        \If{$|\mathcal{D}|\geq B$}{
            Sample mini-batch of size $B$ from $\mathcal{D}$\;
            Minimise $\mathcal{L}(\theta)$ (Eq.~\eqref{eq:td_loss}) via Adam\;
        }
        \If{$f(s_{t+1})=y_\text{target}$ \textbf{and} $r_t>r_\text{best}$}{
            $\bx_\text{best}\leftarrow s_{t+1}$;\quad $r_\text{best}\leftarrow r_t$\;
        }
        \If{$f(s_{t+1})=y_\text{target}$ \textbf{or} step $= T_\text{max}$}{\textbf{break}\;}
    }
    \If{episode $\bmod C = 0$}{$\theta^-\leftarrow\theta$\;}
    Decay $\epsilon_t$ via Eq.~\eqref{eq:eps_decay}\;
}
\Return $\bx_\text{best}$\;
\end{algorithm}

\subsection{Image Explainability Methods}

\subsubsection{Grad-CAM}

For a convolutional network, the class-discriminative localisation map for
class $c$ is
\begin{equation}
L^c_\text{Grad-CAM} = \text{ReLU}\!\left(\sum_k \alpha_k^c A^k\right),
\label{eq:gradcam}
\end{equation}
where $A^k$ is the activation map of the $k$-th feature map in the target
convolutional layer and
\begin{equation}
\alpha_k^c = \frac{1}{Z}\sum_i\sum_j
  \frac{\partial y^c}{\partial A^k_{ij}}
\label{eq:gradcam_weights}
\end{equation}
is obtained by global average pooling of the gradients of the class score $y^c$
with respect to $A^k$. The ReLU in Eq.~\eqref{eq:gradcam} retains only
features that have a positive influence on the target class.

\subsubsection{LIME for Images}

The image is first segmented into $m$ superpixels
$\mathcal{I} = \{S_1,\ldots,S_m\}$ using SLIC or Quickshift. A binary vector
$\mathbf{z}\in\{0,1\}^m$ indicates which superpixels are visible:
$\mathcal{I}'(\mathbf{z})=\bigcup_{i:z_i=1}S_i$. The surrogate is fitted as
\begin{equation}
\bw^* = \argmin_{\bw}\sum_{i=1}^{N}\pi(\mathbf{z}_i)
  \bigl(f(\mathcal{I}'(\mathbf{z}_i)) - \bw^\top\mathbf{z}_i\bigr)^2
  + \lambda\|\bw\|_2^2,
\label{eq:lime_image}
\end{equation}
where $\pi(\mathbf{z})$ is an exponential kernel on the number of masked
superpixels.

%% ─────────────────────────────────────────────────────────────────────────────
\section{Experimental Evaluation}
\label{sec:experiments}
%% ─────────────────────────────────────────────────────────────────────────────

\subsection{Experimental Setup}

\paragraph{Datasets} We evaluate on four tabular datasets---Iris (150 samples,
4 features), UCI Diabetes/Pima (768 samples, 8 features), Credit Scoring
(30,000 samples, 23 features), and Adult Income (48,842 samples, 14
features)---and two image datasets: an ImageNet-1k subset and a chest X-ray
collection (5,000 images).

\paragraph{Models} We evaluate over Random Forest (100 trees), Gradient
Boosting (100 estimators), MLP (128-64-32 hidden units), and ResNet-50 (images).

\paragraph{Baselines} We compare RL-CF against: Random Search, Gradient-Free
optimisation (Nelder-Mead), DiCE~\cite{mothilal2020dice}, and a Genetic
Algorithm (GA) counterfactual solver.

\paragraph{Metrics} We report:
\begin{itemize}
  \item \textbf{Success Rate} (SR): fraction of queries that find a valid CF.
  \item \textbf{Proximity}: $\|\bx-\bxp\|_2$.
  \item \textbf{Sparsity}: $\|\bx-\bxp\|_0$ (number of changed features).
  \item \textbf{Time (s)}: wall-clock seconds per query.
  \item \textbf{Diversity}: mean pairwise $\ell_2$ distance among $K$ returned CFs,
    $\text{Div}=\frac{1}{K(K-1)}\sum_{i\neq j}\|\bx_i'-\bx_j'\|_2$.
\end{itemize}
All results report mean $\pm$ std over 5 independent runs with fixed random seeds.

\subsection{Counterfactual Quality Comparison}

Table~\ref{tab:cf_comparison} reports counterfactual quality on the Credit
Scoring dataset with a Gradient Boosting classifier. RL-CF achieves the highest
success rate and the best proximity and sparsity simultaneously, while also
being the fastest at inference time (after the initial training phase).

\begin{table}[t]
\centering
\caption{Counterfactual Quality on Credit Scoring (GB Classifier)}
\label{tab:cf_comparison}
\begin{tabular}{lcccc}
\toprule
\textbf{Method} & \textbf{SR}$\uparrow$ & \textbf{Prox}$\downarrow$
  & \textbf{Sparse}$\downarrow$ & \textbf{Time (s)}$\downarrow$ \\
\midrule
Random Search   & 0.72 & 3.45 & 8.2 & 12.3 \\
Gradient-Free   & 0.85 & 2.89 & 6.5 &  8.7 \\
DiCE            & 0.91 & 2.34 & 5.8 & 15.2 \\
Genetic Alg.    & 0.88 & 2.67 & 6.1 & 22.5 \\
\textbf{RL-CF (Ours)} & \textbf{0.94} & \textbf{2.12}
  & \textbf{4.9} & \textbf{6.8} \\
\bottomrule
\end{tabular}
\end{table}

\subsection{Ablation Study: Reward Components}

Table~\ref{tab:ablation} progressively adds reward terms to isolate each
component's contribution. The full reward achieves the best performance across
all metrics; removing even the realism term degrades sparsity by 14\%.

\begin{table}[t]
\centering
\caption{Ablation Study --- Impact of Reward Components}
\label{tab:ablation}
\begin{tabular}{lccc}
\toprule
\textbf{Reward Configuration} & \textbf{SR}$\uparrow$
  & \textbf{Prox}$\downarrow$ & \textbf{Sparse}$\downarrow$ \\
\midrule
Classification only           & 0.89 & 3.21 & 9.4 \\
$+$ Proximity                 & 0.92 & 2.45 & 8.1 \\
$+$ Sparsity                  & 0.93 & 2.28 & 5.6 \\
$+$ Realism (Full)            & \textbf{0.94} & \textbf{2.12} & \textbf{4.9} \\
\bottomrule
\end{tabular}
\end{table}

\subsection{Case Study: Credit Scoring}

A loan applicant is rejected ($f(\bx)=0$) with features: income \$35,000,
credit score 620, debt-to-income 0.45, employment years 2. RL-CF recommends
raising the credit score to 680 (+60 points) and reducing debt-to-income to
0.35 ($-0.10$), flipping the prediction to approval with only 2 feature changes.
The traditional counterfactual suggests 8 feature changes to achieve the same
outcome. SHAP identifies credit score as the most influential feature
($\phi=-0.23$) but provides no actionable magnitude---a gap that RL-CF closes.

%% ─────────────────────────────────────────────────────────────────────────────
\section{Discussion}
\label{sec:discussion}
%% ─────────────────────────────────────────────────────────────────────────────

\subsection{Advantages of the RL Approach}

The RL formulation confers several benefits over one-shot optimisation. The DQN
agent \emph{learns from experience}: as it accumulates episodes it refines its
policy, producing sparser and more realistic counterfactuals. The multi-objective
reward naturally balances competing desiderata---class flip, proximity, sparsity,
and distributional realism---without requiring a separate Pareto search. The
epsilon-greedy strategy promotes diversity through exploration, partially
addressing the diversity limitation of single-objective optimisers. Finally, the
trained agent can be \emph{re-used} across queries on the same model, amortising
training cost over multiple explanations.

\subsection{Limitations and Future Work}

Several limitations warrant attention. (i) The discrete action space uses a fixed
step size $\alpha$; a continuous action space (via policy gradient methods such as
PPO or SAC) would be more expressive. (ii) High-dimensional inputs ($d > 100$)
may require hierarchical or attention-based architectures to keep $|\mathcal{A}|
= 2d$ manageable. (iii) Reward weights $\lambda_1,\lambda_2,\lambda_3$ require
domain-specific tuning. (iv) The agent is trained per black-box model; a
meta-learning approach could enable zero-shot transfer.

Future work will explore policy gradient methods for continuous action spaces,
multi-agent RL for explicitly diverse counterfactual generation, integration
with causal inference to ensure counterfactual \emph{validity} beyond statistical
proximity, and formal convergence bounds for the RL-CF policy.

\subsection{Broader Impact}

\emph{expliRL} can support fair lending (actionable feedback to loan applicants),
clinical decision support (interpretable recommendations to clinicians), fraud
detection (explainable flags for analysts), and regulatory compliance (GDPR
right to explanation). Ethical risks include adversarial gaming of the
explanation interface, potential leakage of sensitive model information through
counterfactuals, and the responsibility gap in validating whether recommended
changes are truly feasible.

%% ─────────────────────────────────────────────────────────────────────────────
\section{Implementation and Availability}
\label{sec:impl}
%% ─────────────────────────────────────────────────────────────────────────────

\emph{expliRL} is implemented in Python 3.7+ using \texttt{scikit-learn},
PyTorch, and Gymnasium. It ships with 15+ unit tests (85\% coverage), a CLI
tool, and a REST API built on FastAPI. The library is installable via:
\begin{verbatim}
  pip install expliRL
\end{verbatim}
and a minimal usage example reads:
\begin{verbatim}
from expliRL import RLCounterfactualExplainer
explainer = RLCounterfactualExplainer(model)
explainer.fit(X_train)
cf = explainer.explain(instance)
\end{verbatim}
Source code, documentation, and reproducibility scripts (fixed seeds, Docker
container) are available at \url{https://github.com/soham164/explirl} under the
MIT licence. The library is also indexed on PyPI at
\url{https://pypi.org/project/expliRL/}.

%% ─────────────────────────────────────────────────────────────────────────────
\section{Conclusion}
\label{sec:conclusion}
%% ─────────────────────────────────────────────────────────────────────────────

We presented \emph{expliRL}, a unified open-source framework for explainable AI
that consolidates SHAP, LIME, optimisation-based counterfactuals, and a novel
RL-driven counterfactual generator under a single API. The core contribution is
casting counterfactual generation as an MDP and solving it with a DQN agent
whose multi-objective reward---combining classification success, proximity,
sparsity, and realism---yields a policy that improves with experience and
generalises across queries. Across four tabular benchmarks, RL-CF achieves a
94\% success rate with 15\% better sparsity and 22\% faster convergence than the
best competing baseline. The framework supports tabular and image data, ships
production-ready tooling, and is publicly available under the MIT licence.

%% ─────────────────────────────────────────────────────────────────────────────
\section*{Acknowledgments}
The authors thank the open-source communities behind \texttt{shap}, \texttt{lime},
\texttt{dice-ml}, and Gymnasium. Computational experiments were conducted on
[hardware platform].

%% ─────────────────────────────────────────────────────────────────────────────
\bibliographystyle{IEEEtran}
\bibliography{references}

%% ─── Inline bibliography (no external .bib required for self-contained file) ─
\begin{thebibliography}{99}

\bibitem{gdpr2016}
European Parliament and Council, ``Regulation (EU) 2016/679 (General Data
Protection Regulation),'' \emph{Official Journal of the European Union},
vol.~L~119, pp.~1--88, 2016.

\bibitem{lundberg2017shap}
S.~M.~Lundberg and S.-I.~Lee, ``A unified approach to interpreting model
predictions,'' in \emph{Advances in Neural Information Processing Systems
(NeurIPS)}, vol.~30, 2017.

\bibitem{ribeiro2016lime}
M.~T.~Ribeiro, S.~Singh, and C.~Guestrin, ```Why should I trust you?'
Explaining the predictions of any classifier,'' in \emph{Proc. ACM SIGKDD
Int. Conf. Knowledge Discovery \& Data Mining (KDD)}, 2016,
pp.~1135--1144.

\bibitem{wachter2017counter}
S.~Wachter, B.~Mittelstadt, and C.~Russell, ``Counterfactual explanations
without opening the black box: Automated decisions and the GDPR,''
\emph{Harvard Journal of Law \& Technology}, vol.~31, no.~2, pp.~841--887,
2018.

\bibitem{mothilal2020dice}
R.~K.~Mothilal, A.~Sharma, and C.~Tan, ``Explaining machine learning
classifiers through diverse counterfactual explanations,'' in \emph{Proc.
ACM Conf. Fairness, Accountability, and Transparency (FAT*)}, 2020,
pp.~607--617.

\bibitem{poyiadzi2020face}
R.~Poyiadzi, K.~Sokol, R.~Santos-Rodriguez, T.~De~Bie, and P.~Flach,
``FACE: Feasible and actionable counterfactual explanations,'' in \emph{Proc.
AAAI/ACM Conf. AI, Ethics, and Society}, 2020, pp.~344--350.

\bibitem{mnih2015dqn}
V.~Mnih et al., ``Human-level control through deep reinforcement learning,''
\emph{Nature}, vol.~518, no.~7540, pp.~529--533, 2015.

\bibitem{selvaraju2017gradcam}
R.~R.~Selvaraju, M.~Cogswell, A.~Das, R.~Vedantam, D.~Parikh, and
D.~Batra, ``Grad-CAM: Visual explanations from deep networks via
gradient-based localization,'' in \emph{Proc. IEEE Int. Conf. Computer
Vision (ICCV)}, 2017, pp.~618--626.

\bibitem{arrieta2020xai}
A.~B.~Arrieta et al., ``Explainable Artificial Intelligence (XAI):
Concepts, taxonomies, opportunities and challenges toward responsible AI,''
\emph{Information Fusion}, vol.~58, pp.~82--115, 2020.

\bibitem{gaudel2010rl}
R.~Gaudel and M.~Sebag, ``Feature selection as a one-player game,'' in
\emph{Proc. Int. Conf. Machine Learning (ICML)}, 2010, pp.~359--366.

\bibitem{guo2021rl}
C.~Guo, F.~Juefei-Xu, Q.~Xie, L.~Ma, L.~Li, W.~Liu, J.~Zhao, and
Y.~Cheng, ``A reinforcement learning-based approach to automatically test
deep neural networks,'' in \emph{Proc. 43rd Int. Conf. Software Engineering
(ICSE)}, 2021, pp.~800--811.

\end{thebibliography}

\end{document}